\documentclass[10pt,a4paper]{article}
\usepackage[margin=18mm]{geometry}
\usepackage{fontspec}
\usepackage{amsmath,amssymb}
\usepackage{booktabs}
\usepackage{tabularx}
\usepackage{enumitem}
\usepackage{xcolor}
\usepackage{hyperref}
\usepackage{listings}

\setmainfont{Noto Sans CJK KR}[AutoFakeSlant=0.2]
\setsansfont{Noto Sans CJK KR}
\setmonofont{D2Coding}
\newfontfamily\latinfont{DejaVu Sans}
\setlength{\parindent}{0pt}
\setlength{\parskip}{5pt}
\setlist{nosep,leftmargin=6mm}
\pagestyle{plain}
\sloppy
\hypersetup{
  colorlinks=true,
  linkcolor=blue,
  urlcolor=blue
}
\lstset{
  basicstyle=\ttfamily\small,
  frame=single,
  columns=fullflexible,
  breaklines=true,
  keepspaces=true
}

\NeedsTeXFormat{LaTeX2e}
\ProvidesPackage{answer_coversheet}[2026/07/30 2026 cryptanalysis contest answer cover sheet]

% Recreates the official 2026_암호분석경진대회-답안양식.hwp look: a
% "2026 암호분석경진대회" / "N번 문제" header (single rule under the first
% row, double rule under the second, matching the HWP's own
% borderfill-4/5/6 cell borders) and a page border, drawn as a background
% layer on every page. Deliberately does NOT wrap body content in any
% width-changing environment (an earlier mdframed-based version looped
% forever once a wide table met its narrowed frame width) so it never
% interacts with a document's existing tables, margins, or line widths.

\RequirePackage{geometry}
\RequirePackage{eso-pic}
\RequirePackage{tikz}
\RequirePackage{fontspec}

% Reserve room at the top of every page for the header below, and enough
% bottom clearance to stay clear of the border (19.3mm inset), without
% touching left/right margins (existing wide tables are sized against those).
\geometry{top=30mm, bottom=22mm}

\newcommand{\ProblemNumber}{}
\newcommand{\SetProblemNumber}[1]{\renewcommand{\ProblemNumber}{#1}}

% Match the font every other submission already uses for its body text
% (the HWP's real font, 함초롬바탕/Hancom HCR Batang, is a proprietary font
% not available via apt/pip here). Override with \renewcommand when the
% document's main font already has no Hangul coverage (see submissions/02).
\IfFontExistsTF{Noto Sans CJK KR}{%
  \newfontfamily\AnswerCoverFont{Noto Sans CJK KR}[AutoFakeSlant=0.2]%
}{%
  \newcommand{\AnswerCoverFont}{}%
}

% Box/row measurements below are taken directly from the HWP's own layout
% (via hwp5html): HeaderPageFooter margin-top 10mm + page padding-top 5mm =
% 15mm top inset; margin-bottom 5mm + padding-bottom 5mm = 10mm, plus the
% content cell's own extra spacing brings it to ~19.3mm; table margin-left
% 1mm inside a 10mm side margin = 11mm; each header row is 6.15mm tall.
% The title/problem-number header (and its rules) is drawn on page 1 only;
% the outer border repeats on every page.
\AddToShipoutPictureBG{%
  \AnswerCoverFont
  \begin{tikzpicture}[overlay]
    \coordinate (TL) at ([xshift=11mm]current page.north west);
    \coordinate (TR) at ([xshift=-11mm]current page.north east);
    \coordinate (BR) at ([xshift=-11mm]current page.south east);
    \draw[line width=0.6pt] ([yshift=-15mm]TL) rectangle ([yshift=19.3mm]BR);
    \ifnum\value{page}=1
      \draw[line width=0.6pt] ([yshift=-21.15mm]TL) -- ([yshift=-21.15mm]TR);
      \draw[line width=0.4pt] ([yshift=-27.1mm]TL) -- ([yshift=-27.1mm]TR);
      \draw[line width=0.4pt] ([yshift=-27.5mm]TL) -- ([yshift=-27.5mm]TR);
      \node[anchor=north] at ([yshift=-15.5mm]current page.north) {2026 암호분석경진대회};
      \node[anchor=north] at ([yshift=-21.65mm]current page.north) {\ProblemNumber 번 문제};
    \fi
  \end{tikzpicture}%
}

\endinput

\SetProblemNumber{8}

\title{블록암호}
\author{중간 상태 누출을 이용한 변형 AES 마스터 키 복원}
\date{2026년 7월}

\begin{document}
\maketitle

\section*{결과 요약}

복원한 16바이트 마스터 키는 다음과 같다.

\begin{center}
\fbox{\texttt{2923be84e16cd6ae529049f1f1bbe9eb}}
\end{center}

문제의 0-based index \texttt{28178}에 해당하는 평문ㆍ암호문과
\texttt{leaked.txt}를 결합했다. 첫ㆍ마지막 라운드 식에서
\texttt{K[0]=29}, \texttt{K[1]=23}, \texttt{K[3]=84}를 직접 구하고,
$X^3$의 $2^{16}$개 완성과 $X^2$ 첫 열의 16개 완성을 공유
마스터 키 바이트로 결합했다. 중복을 제거한 full key 중 누출 대상
암호문을 만족하는 키는 정확히 하나였으며, 그 키로 제공된 50,000개
평문을 모두 다시 암호화한 결과 불일치는 0개였다.

\begin{center}
\begin{tabular}{lr}
\toprule
검증 항목 & 결과 \\
\midrule
누출 대상 index & \texttt{28178} (0부터 시작) \\
직접 복원한 마스터 키 바이트 & \texttt{K[0], K[1], K[3]} \\
누출 대상 pair를 만족하는 서로 다른 full key & 1개 \\
전체 재암호화 & 50,000개 중 mismatch 0개 \\
\bottomrule
\end{tabular}
\end{center}

\section{문제 암호와 상태 표현}

AES state는 표준의 column-major 순서를 사용한다.
\[
  \mathit{index}=\mathit{row}+4\mathbin{\cdot}\mathit{column}.
\]
SubBytes, ShiftRows, MixColumns를 합쳐
\[
  F=\operatorname{MC}\circ\operatorname{SR}\circ\operatorname{SB}
\]
라 쓰면 문제의 변형 7라운드 암호는
\[
  X^0=P,\qquad X^{i+1}=F(X^i)\oplus k_i\quad(0\le i<7),
  \qquad C=X^7
\]
이다. 표준 AES와 달리 초기 AddRoundKey가 없고 마지막 라운드에도
MixColumns가 있다. 또한 라운드 키는 표준 AES key expansion이 아니라,
마스터 키 바이트의 문제 지정 permutation에 일부 라운드의 MixColumns를
적용해 만든다. 특히 $k_0=K$이며, 서로 떨어진 라운드가 같은 마스터 키
바이트를 공유한다. 이 공유 관계가 여러 중간 상태 누출을 작은
key-byte 제약으로 연결하게 해 준다.

\subsection{힌트와 실제 공격 범위}

문제는 Demirci--Selçuk meet-in-the-middle 공격을 힌트로 준다. 원 논문은
5라운드 distinguisher를 바탕으로 reduced-round AES를 공격한다. 그러나
본 solver는 그 distinguisher, 논문의 다라운드 backward 경로, 또는
논문에 제시된 AES-192/256 공격 복잡도를 구현하지 않는다.

여기서는 문제가 훨씬 강한 중간 상태 bit 누출과 특수한 라운드 키
permutation을 직접 제공한다. 원 공격에서 가져온 것은 중간 표현의
후보를 미리 계산하고 공통 값으로 결합한다는 관점이며, 실제 구현은 이
누출을 이용해 앞뒤 방향에서 key-byte 제약을 직접 결합하는 방식이다.

\section{누출 대상과 직접 복원}

\subsection{0-based index 검증}

배열의 \texttt{28178}번째 block을 읽은 결과는 다음과 같다.

\begin{lstlisting}
P[28178] = b3a6db3c870c3e99245e0d1c06b747de
C[28178] = 0b764cbfd5a347a13b41b4fdd5b11e00
\end{lstlisting}

두 값은 \texttt{leaked.txt}의 $X^0$, $X^7$ 전체 byte와 일치한다.
따라서 자연어의 ``28,179번째''를 1-based 배열 첨자로 오해하지 않고
문제에 명시된 0-based index를 사용했다.

\subsection{첫 라운드의 \texttt{K[0]}}

$k_0=K$이고 $X^1=F(X^0)\oplus k_0$이다. 누출에는
$X^1[0]=\texttt{20}$이 완전히 주어지므로
\[
 K[0]=F(X^0)[0]\oplus X^1[0]=\texttt{29}
\]
를 바로 얻는다.

\subsection{마지막 라운드의 \texttt{K[1]}, \texttt{K[3]}}

마지막 라운드 키는
\[
k_6=\operatorname{MC}\bigl(
 K[1],K[5],K[7],K[14],\ldots,
 K[9],K[13],K[3],K[6],\ldots\bigr)
\]
형태다. $X^7=F(X^6)\oplus k_6$의 각 열에 InvMixColumns를 적용하면
MixColumns의 선형성 때문에 다음 관계가 된다.
\[
\begin{aligned}
\operatorname{InvMC}(C_{\mathrm{col}\,0})
  \oplus [K_1,K_5,K_7,K_{14}]
  &=\operatorname{SR}(\operatorname{SB}(X^6))_{\mathrm{col}\,0},\\
\operatorname{InvMC}(C_{\mathrm{col}\,2})
  \oplus [K_9,K_{13},K_3,K_6]
  &=\operatorname{SR}(\operatorname{SB}(X^6))_{\mathrm{col}\,2}.
\end{aligned}
\]
$X^6[0]$과 $X^6[2]$가 완전히 누출되어 있으므로
\[
 K[1]=\texttt{23},\qquad K[3]=\texttt{84}
\]
를 직접 복원한다. 이로써 탐색 전에 마스터 키 3바이트가 확정된다.

\section{누출 completion과 후보 결합}

최종 구현에서 사용하는 누출은 다음과 같다.

\begin{itemize}
  \item $X^1[0]$: 완전한 1바이트;
  \item $X^2[0..3]$: 각 바이트에 unknown bit 1개;
  \item $X^3[0..15]$: 각 바이트에 unknown bit 1개;
  \item $X^4[0],X^4[5],X^4[10],X^4[15]$: 각 바이트에 unknown bit 1개;
  \item $X^6[0],X^6[2]$: 완전한 2바이트; 그리고
  \item $X^0$, $X^7$: 완전한 평문과 암호문.
\end{itemize}

\subsection{\texorpdfstring{$X^3\rightarrow X^4$}{X3 to X4} 대각선}

$X^3$의 각 바이트에는 unknown bit가 하나이므로 가능한 전체 state는
$2^{16}=65{,}536$개다. 각 completion마다 $F(X^3)$와 필요한 네
InvMixColumns 열을 미리 계산해 record로 보관한다. $X^4$의 대각 위치
$0,5,10,15$와 $k_3$의 단순 마스터 키 permutation을 대조하면 다음
후보 집합을 얻는다.

\begin{center}
\begin{tabular}{lr}
\toprule
마스터 키 바이트 & 서로 다른 후보 수 \\
\midrule
\texttt{K[10]} & 32 \\
\texttt{K[11]} & 16 \\
\texttt{K[12]} & 32 \\
\texttt{K[13]} & 32 \\
\bottomrule
\end{tabular}
\end{center}

\subsection{\texorpdfstring{$X^2$}{X2} 첫 열과 compact relation}

$X^2[0..3]$에도 각각 unknown bit가 하나뿐이므로 첫 열 completion은
$2^4=16$개다. MixColumns의 선형성을 이용하면
\[
\operatorname{InvMC}(X^2_{\mathrm{col}\,0})
=\operatorname{SB}(X^1_{\mathrm{diag}})
 \oplus [K_5,K_8,K_6,K_3]
\]
의 형태로 바꿀 수 있다. 이미 확정된 \texttt{K[3]}과
\texttt{K[10]} 후보 32개를 결합하면
\[
 16\mathbin{\cdot}32=512
\]
개의 기본 $X^2$/key relation만 남는다. 이 단계에서 서로 다른 후보
수는 \texttt{K[5]}, \texttt{K[8]}, \texttt{K[15]}가 각각 16개이고,
\texttt{K[6]}은 192개다.

\subsection{실제 leak-assisted join}

solver의 결합 순서는 다음과 같다.

\begin{enumerate}
  \item 앞뒤 직접식으로 \texttt{K[0]}, \texttt{K[1]},
        \texttt{K[3]}을 확정한다.
  \item $X^3$의 $2^{16}$개 completion을 record로 만든다.
  \item 각 record의 $F(X^3)$ 대각 byte와 $X^4$ 대각 누출을 결합해
        \texttt{K[10..13]} 후보를 만든다.
  \item $X^2$의 16개 completion과 $X^3$ record를 공유 key byte로
        비교한다. 이 직접 비교 축은
        $16\mathbin{\cdot}2^{16}=2^{20}=1{,}048{,}576$개다.
  \item 살아남은 partial key에서 \texttt{K[2]}, \texttt{K[4]},
        \texttt{K[14]}를 전방ㆍ역방향 $X^2$가 완전히 같다는 조건으로
        채운다.
\end{enumerate}

최종 코드는 명시적인 $X^5$ 역계산이나 $X^7\rightarrow X^4$ backward
table을 만들지 않는다. 여러 라운드 키에 재사용된 동일 마스터 키
바이트가 곧 join key이며, 누출이 그 값들을 빠르게 고정한다.

\section{마지막 3바이트의 완전 탐색과 유일성}

\texttt{solve\_aes\_key.py} 내의 \texttt{complete\_partial\_keys()}는
한 partial key에서 처음 발견한
조합을 반환하지 않는다. 먼저 \texttt{K[2]}와 \texttt{K[4]}를 각각
0부터 255까지 별도 목표 byte로 거른 뒤, 살아남은 모든 조합에 대해
\texttt{K[14]}의 256개 값을 끝까지 조사한다. 후보마다
\[
 X^2_{\mathrm{forward}}(K)
 =F(X^1)\oplus k_1
\]
와
\[
 X^2_{\mathrm{backward}}(K)
 =F^{-1}(X^3\oplus k_2)
\]
의 16바이트 전체가 같은지 확인한다.

생성된 모든 full key는 \texttt{set[bytes]}에 넣어 서로 다른 경로가
같은 키를 만드는 경우를 중복 제거한다. 그 뒤 누출 대상 평문을 직접
암호화해 해당 암호문과 일치하는 키만 남기고, 서로 다른 생존 key가
정확히 하나임을 확인한 뒤에 결과를 반환한다.

\section{복잡도와 측정 범위}

\begin{center}
\begin{tabularx}{\textwidth}{Xr}
\toprule
단계 & 후보 또는 작업 수 \\
\midrule
$X^3$의 16개 unknown bit completion & $2^{16}=65{,}536$ \\
$X^2[0..3]$의 4개 unknown bit completion & $2^4=16$ \\
기본 $X^2$/\texttt{K[10]} relation & $16\mathbin{\cdot}32=512$ \\
$X^2$ completion과 $X^3$ record 비교 & $2^{20}=1{,}048{,}576$ \\
최종 사후 검증 & 50,000 blocks $\mathbin{\cdot}$ 7 rounds \\
\bottomrule
\end{tabularx}
\end{center}

메모리의 주된 항은 65,536개 \texttt{x3\_records}다. 현재 Python
object 표현에서는 공간 복잡도가 $O(2^{16})$이며, packed integer나
고정폭 array를 쓰면 상수 계수는 줄일 수 있다.

Python 3.11.2, Linux 6.1.0-51-cloud-amd64, glibc 2.36 환경에서 key
탐색, 완전 completion, 50,000쌍 검증을 함께 잰 두 번의 실행은
41.706--57.160초, peak RSS 59,812--60,436 KiB(약 58.4--59.0 MiB)였다.
다회 benchmark가 아니므로 환경에 따라 달라질 수 있는 참고값이다.

탐색 초기에는 \texttt{z3-solver}로 full AES bit-vector 모델을 짧게
시도했으나 timeout으로 끝나 포기했다. 최종적으로 leak-assisted join을
택한 것은 이 방법이 실제 key를 복원하고 50,000쌍 전체로 검증됐기
때문이다.

\section{구현, 재현 및 검증}

최종 Python solver는 AES 정ㆍ역연산, 문제의 라운드 키 생성, leak
parser, 후보 record와 join을 구현한다. 마지막 세 key byte를 빠짐없이
생성하고 전체 데이터를 검증하는 과정도 Python 표준 라이브러리만
사용한다.

\begin{lstlisting}
python3 investigate_aes_leak.py
python3 solve_aes_key.py
\end{lstlisting}

최종 solver의 핵심 출력은 다음과 같다.

\begin{lstlisting}
master key = 2923be84e16cd6ae529049f1f1bbe9eb
verified pairs = 50000
mismatches = 0
submission key check = PASS
\end{lstlisting}

실행 중에는 다음도 함께 검사한다.

\begin{itemize}
  \item 누출 index의 $X^0$, $X^7$과 제공 파일의 평문ㆍ암호문 일치;
  \item 직접 복원한 3바이트와 누출 bit pattern의 일치;
  \item 모든 completion에서 생성한 full key의 중복 제거와 유일성;
  \item 복원 key로 50,000개 block 전체를 재암호화한 mismatch 0; 그리고
  \item 출력 key와 \texttt{master\_key.txt}의
        byte-for-byte 일치.
\end{itemize}

\section{참고 자료}

\begin{itemize}
  \item \href{https://doi.org/10.6028/NIST.FIPS.197-upd1}
        {NIST, \emph{FIPS 197: Advanced Encryption Standard (AES)},
        2023 update}. Column-major state, SubBytes, ShiftRows, MixColumns와
        역변환 구현의 기준으로 사용했다.
  \item \href{https://www.iacr.org/archive/fse2008/50860106/50860106.pdf}
        {{\latinfont Hüseyin Demirci and Ali Aydın Selçuk},
        \emph{A Meet-in-the-Middle Attack on 8-Round AES}, FSE 2008}.
        중간 표현 후보의 precompute/join 관점을 참고했으며, 본 solver가
        논문의 distinguisher와 공격 전체를 구현하지 않았음을 본문에서
        구분했다.
  \item \path{8_블록암호.zip}. 변형 7라운드 구조, 라운드 키
        permutation, 50,000쌍과 중간 상태 누출의 원문 자료다.
  \item \path{solve_aes_key.py}. 누출 completion,
        key-byte constraint join, exhaustive full-key 생성과 전체 검증을
        재현하는 구현이다.
\end{itemize}

\end{document}
